%%%%%%%%%%%%%%%%%%%%%%%%%%%%%%%%%%%%%%%%%%%%%%%%%%%%%%%%%%%%%%%%%%%%%%%%%%%%%%%%%%%%%%%%%%%%%%%%%%%%%%%%%
%%%%%  preamble.tex -- shared settings for every Module 1 topic file                                %%%%%
%%%%%  Loaded by main.tex and, through subfiles, by each topic in topics/                           %%%%%
%%%%%%%%%%%%%%%%%%%%%%%%%%%%%%%%%%%%%%%%%%%%%%%%%%%%%%%%%%%%%%%%%%%%%%%%%%%%%%%%%%%%%%%%%%%%%%%%%%%%%%%%%

\usepackage[a4paper,margin=2.54cm]{geometry}
\usepackage[T1]{fontenc}
\usepackage{mathptmx}                 
\usepackage[scaled=0.92]{helvet}      
\usepackage{microtype}
\usepackage{setspace}
\setstretch{1.2}
\setlength{\parindent}{0pt}
\setlength{\parskip}{6pt}

\usepackage{array,tabularx,booktabs}
\renewcommand{\arraystretch}{1.25}
\newcolumntype{L}{>{\raggedright\arraybackslash}X}
\usepackage{enumitem}
\setlist{itemsep=2pt,topsep=4pt,leftmargin=*}
\usepackage[font=small,labelfont=bf,textfont=it,skip=6pt]{caption}

\usepackage{titlesec}
\titleformat{\section}{\normalfont\large\bfseries}{\thesection.}{0.5em}{}
\titleformat{\subsection}{\normalfont\normalsize\bfseries}{\thesubsection}{0.6em}{}
\titlespacing*{\section}{0pt}{14pt}{6pt}
\titlespacing*{\subsection}{0pt}{10pt}{4pt}

\usepackage{fancyhdr}
\pagestyle{fancy}
\fancyhf{}
\fancyhead[L]{\small 26ESCC106 \,|\, Module 1}
\fancyhead[R]{\small\rightmark}
\fancyfoot[C]{\thepage}
\renewcommand{\headrulewidth}{0.4pt}
\renewcommand{\sectionmark}[1]{}      
\renewcommand{\subsectionmark}[1]{}  

\usepackage[most]{tcolorbox}
\newtcolorbox{notebox}{colback=white,colframe=black,boxrule=0.6pt,
  arc=0pt,left=8pt,right=8pt,top=6pt,bottom=6pt,fontupper=\itshape}

\usepackage{tikz}
\usetikzlibrary{arrows.meta,calc}
\tikzset{
  every picture/.style={line width=0.8pt,>={Stealth[length=2.2mm]}},
  lbl/.style={font=\sffamily\footnotesize,align=center},
}
\NewDocumentCommand{\B}{O{} O{} m m m}{%
  \draw[#1] (#3) rectangle (#4);
  \path (#3) -- node[lbl,#2]{#5} (#4);}
\newcommand{\monitor}[2]{%
  \draw[line width=2.2pt,fill=white] (#1-0.55,#2-0.35) rectangle (#1+0.55,#2+0.37);
  \fill (#1-0.10,#2-0.35) -- (#1+0.10,#2-0.35) -- (#1+0.13,#2-0.55) -- (#1-0.13,#2-0.55) -- cycle;
  \fill (#1-0.35,#2-0.62) rectangle (#1+0.35,#2-0.54);}

\usepackage[hidelinks]{hyperref}
\urlstyle{same}
\newcommand{\rc}[1]{[\ref{#1}]}
\newcommand{\rcc}[2]{[\ref{#1},\,\ref{#2}]}

%%%%%%%%%%%%%%%%%%%%%%%%%%%%%%%%%%%%%%%%%%%%%%%%%%%%%%%%%%%%%%%%%%%%%%%%%%%%%%%%%%%%%%%%%%

%%%%%  Topic heading %%%%%
%  \topic{<name>} starts a new topic: it prints the heading, puts the name in
%  the running header, restarts section numbering at 1 and numbers the figures
%  of that topic <topic no.>.<figure no.>.
\newcounter{topicno}
\newcommand{\topicname}{}
\renewcommand{\thefigure}{\arabic{topicno}.\arabic{figure}}
\newcommand{\topic}[1]{%
  \ifnum\value{topicno}>0 \clearpage\fi
  \stepcounter{topicno}%
  \renewcommand{\topicname}{#1}%
  \markright{#1}%
  \setcounter{section}{0}\setcounter{figure}{0}%
  \begin{center}{\Large\bfseries Topic in Module 1 --- #1}\end{center}}
